% ======================================================================
%  Section 4.4.3 --- The Hermite-informed weight.  Corrected.
%
%  THREE SUBSTANTIVE PROBLEMS.
%
%  A. EQUATION (4.11) NORMALISES BY THE WRONG COUNT. It divides by P, the
%     number of measurements. G is defined per image POSITION, and there
%     are four measurements per position, so the mean must divide by the
%     number of positions. The implementation agrees:
%         double Gmean = sumG / nbIn;    // nbIn = 66000, not 264000
%     As written the thesis mean is four times too small, hence Gbar four
%     times too large, and the floor at 0.1 would bite quite differently.
%     Fixed by introducing N for the number of positions.
%
%  B. "the occluded region is rendered as a uniform DARK area ... so its
%     Hermite response energy falls well below the image average". With the
%     DC-free measure of 4.4.2, darkness is irrelevant: G has exactly zero
%     DC gain, so a uniform BRIGHT occluder gives the same vanishing
%     response. Correcting this strengthens the section, because it leaves
%     the textured occluder as the only case the mechanism misses -- which
%     is exactly what 4.4.4 says.
%
%  C. "its measurements are among the FIRST to be rejected" contradicts
%     4.7.3, which measures the opposite: at iteration 1 only 4.2% of
%     measurements are rejected against 7.4% covered by the disc, so the
%     occluder is not yet being rejected at all. It is rejected decisively
%     at convergence, not early. The two sections cannot both stand.
%
%  Plus one addition that answers an obvious objection, and the usual
%  labels and punctuation.
% ======================================================================

\subsection{The Hermite-informed weight}
\label{sec:occlusion-mechanism}
%%% 4.4.2 refers to this subsection for the vanishing-response property.

The structure measure is used to normalise the residual before it is submitted
to the robust estimator. It is first expressed relative to its mean over the
image,
%%% (A) N, not P. G lives on positions; P counts measurements, of which
%%% there are four per position.
\begin{equation}
    \bar{G}(p)=\frac{G(p)}{\overline{G}},
    \qquad
    \overline{G}=\frac{1}{N}\sum_{p}G(p),
    \label{eq:Gbar}
\end{equation}
where $N$ is the number of image positions carrying a measurement, so that
$P=4N$ for the four-scale feature vector of Chapter~\ref{ch:herpvs}. The
measure is thereby dimensionless and self-calibrating: it reports the
structure at $p$ relative to the average structure of the current image, and
is therefore insensitive to the overall texture level of the scene. The
normalised ratio is floored,
\begin{equation}
    \bar{G}(p)\leftarrow\max\bigl(\bar{G}(p),\,\bar{G}_{\min}\bigr),
    \qquad \bar{G}_{\min}=0.1,
    \label{eq:Gfloor}
\end{equation}
and the structure-normalised residual is
\begin{equation}
    \tilde{e}_{j}=\frac{e_{j}}{\bar{G}(p_{j})},
    \label{eq:normalized-residual}
\end{equation}
where $p_{j}$ is the image position of the $j$-th measurement, so that the
four measurements taken at a given position share a single scale factor.

The floor bounds the inflation applied to any residual at
$1/\bar{G}_{\min}=10$, so that a measurement in a structureless region is
penalised strongly but not without limit: such a measurement is rejected if
its residual remains large after a bounded rescaling, not unconditionally.
%%% Sharpened. "caps the rejection strength" suggested the floor prevents
%%% rejection, which it does not; it prevents rejection from becoming
%%% automatic.
This bound also matters for the rank condition of
Section~\ref{sec:weighted-law}: an unbounded inflation would drive every
exactly flat position to zero weight, and a sufficiently untextured scene
could then remove a direction of $\mathbf{H}_{\mathbf{D}}$ altogether.
%%% Reference was sec:robust-framework, the whole of 4.2; the rank
%%% condition and the floor that guards it are in 4.2.2.

The weights are then obtained by applying the Tukey function~\eqref{eq:tukey}
to the structure-normalised residuals,
\begin{equation}
    w^{\mathrm{H}}_{j}=w^{\mathrm{T}}\!\left(\frac{\tilde{e}_{j}}{\hat{s}_{\tilde{e}}}\right),
    \qquad
    \hat{s}_{\tilde{e}}=1.4826\operatorname{med}_{j}
    \bigl\lvert\tilde{e}_{j}-\operatorname{med}_{k}(\tilde{e}_{k})\bigr\rvert,
    \label{eq:hermite-weights}
\end{equation}
and assembled into
$\mathbf{D}^{\mathrm{H}}=\mathrm{diag}(w^{\mathrm{H}}_{1},\dots,w^{\mathrm{H}}_{P})$.
Here $\hat{s}_{\tilde{e}}$ is the robust scale of the
\emph{structure-normalised} residuals, computed as in
Section~\ref{sec:mad} but on $\tilde{e}_{j}$ rather than on $e_{j}$; the
normalisation therefore precedes the scale estimation, and the two variants
differ in the argument supplied to the estimator rather than in the estimator
itself.

%%% ADDED. This answers the first objection a reader will raise: if the
%%% scale is re-estimated on the normalised residuals, does the
%%% normalisation not simply cancel? The answer is instructive and takes
%%% two sentences.
That the scale is re-estimated on $\tilde{e}$ rather than carried over from
$e$ has a consequence worth stating. Were $\bar{G}$ constant over the image,
the normalisation would multiply every residual by the same factor, that
factor would appear identically in $\hat{s}_{\tilde{e}}$, and the ratio
$\tilde{e}_{j}/\hat{s}_{\tilde{e}}$ would be unchanged: a uniform rescaling is
inert. What acts is therefore not the magnitude of the normalisation but its
\emph{variation} across the image, which redistributes the residuals relative
to their common scale and so changes which of them are extreme. The
weighting is, in this sense, a reallocation of the outlier budget rather than
a change in its size.

The interpretation is a rescaling of the outlier threshold according to the
local image content. Where the structure is stronger than the image average,
$\bar{G}>1$, the normalised residual is reduced, and a measurement must depart
substantially further from the model before it is rejected: large residuals
are expected there and are treated as informative. Where the structure is
weaker than average, $\bar{G}<1$, the normalised residual is inflated, and the
measurement is rejected on comparatively slight evidence: little pose
information is available there and the measurement is not worth defending. The
estimator thus retains exactly the redescending behaviour of
Section~\ref{sec:tukey-biweight}, but applies it to a residual that has been
placed on a common footing across regions of differing image structure.
%%% The double comma after "reduced,," is gone.

%%% (B) and (C). "dark" removed; "among the first to be rejected" corrected
%%% to what 4.7.3 measures.
In the configurations considered here the occluded region is a uniform area of
the target surface. Since $G$ is built from zero-mean kernels its response
vanishes there whatever the intensity of that area, so $\bar{G}$ meets the
floor~\eqref{eq:Gfloor}, the residual is inflated tenfold, and the
measurement is rejected. It is the uniformity of the region and not its
brightness that produces this: a uniform bright occluder would be suppressed
identically. This is the mechanism by which the weighting acts on occlusion,
and Section~\ref{sec:weight-maps} shows it operating. It should be said that
the rejection is not immediate. While the pose error is large the residual is
large over the whole image and the robust scale
$\hat{s}_{\tilde{e}}$ is large in proportion, so the occluded measurements do
not yet stand out; they are isolated as the pose converges and the scale
falls, which Section~\ref{sec:weight-maps} measures. A strongly textured
occluder, by contrast, would produce a large $G$ at every stage and would not
be suppressed by this mechanism at all; the scope of the method is discussed
in Section~\ref{sec:hermite-weight-limitation}.

% ======================================================================
%  LABELS: SETTLE 4.4.4 NOW. Three different labels for it are in play --
%  sec:hermite-weight-limitation (yours, used above), sec:ablation-measure
%  and sec:acknowledged (both mine, from earlier drafts of 4.4.2). Keep
%  YOURS and change the reference in 4.4.2 to match. 4.4.4 will host both
%  the acknowledged limitation and the ablation of the measure, so one
%  label serves.
%
%  Other references corrected above:
%    rank condition   sec:robust-framework -> sec:weighted-law   (4.2.2)
%    robust scale     sec:tukey            -> sec:mad            (4.3.2)
%    redescending     sec:tukey            -> sec:tukey-biweight (4.3.1)
%
%  NOTATION: D^H is the fourth weight matrix written with the same bold D
%  used for the Gauss-Hermite kernels of 4.4.2. See the note in the 4.4.2
%  correction; if you adopt W for the weight matrices, this becomes W^H.
%
%  CONSEQUENT EDIT ELSEWHERE
%  4.4.4 currently opens "The structure measure (4.10) reports that the
%  image is locally flat, not why it is flat." That is right and is now the
%  natural continuation of the paragraph above. But its next sentence says
%  a flat region is down-weighted "whatever the cause of its flatness",
%  which is the same point as the brightness clarification here -- check
%  that the two do not now repeat each other.
% ======================================================================
